\documentclass[11pt,a4paper]{article}
\usepackage{amsmath,booktabs,geometry}
\geometry{margin=2.3cm}
\setlength{\parskip}{0.35em}
\setlength{\parindent}{0em}
\title{Which feature layer carries a drifted Tree-of-Thoughts}
\author{}
\date{2026}
\begin{document}
\maketitle

\begin{abstract}
FSDS is applied to two searches after the features are grouped into six layers.
On Game of 24 the outcome rate falls from \(1.00\) to \(0.25\), and the score-shape layers take \(65\%\) of the concept-drift SSE.
The graph layer takes \(9\%\).
On the one-step offer search the tree never changes shape, the graph layer takes nothing, and the previous outcome takes \(73\%\).
The four-number state encoding moves sharply and still fails the concept test, because it is the same step as the offer change.
\end{abstract}

\section{What is extracted}
Each finished episode supplies one row.
\(Y\) is \(1\) when the search returns a solved state, and \(0\) otherwise.
The current episode's residual is not a column.
The layers used here are a small subset of a six-layer scheme, with sizes fixed by the search rather than by a 768-dimensional text encoder.

\begin{center}
\begin{tabular}{ll}
\toprule
Layer & Columns on the committed path \\
\midrule
L1 graph & depth, children, siblings, path length, branch count \\
L2 history & previous score, score change, rolling mean and std, previous \(Y\) \\
L3 statistics & score std, score \(z\), score entropy \\
L4 embedding & four numbers: sorted state entries divided by 24, or a 4-bin identity \\
L5 interaction & children per depth, depth times score gap, score times path length \\
L6 context & beam, maximum depth, last-step fraction \\
\bottomrule
\end{tabular}
\end{center}

Concept drift for a column \(x\) is the SSE drop from \(Y\sim 1+\mathrm{batch}\) to \(Y\sim 1+\mathrm{batch}+x+\mathrm{batch}\cdot x\), split after 12 episodes, with an 80-shuffle \(p\)-value.
A second test records the absolute mean shift across the drift at episode 16.
Layer share is the sum of positive SSE drops inside the layer, divided by the sum over layers.

\section{Game of 24}
Beam 4, depth 3, drift at episode 16, judge mix \(0.28\).
Forty episodes.
\(Y\) moves from \(1.00\) to \(0.25\).

\begin{center}
\begin{tabular}{lrrrl}
\toprule
Layer & SSE share & Best column & Best \(p\) & Largest mean shift \\
\midrule
L3 statistics & 0.354 & score entropy & 0.012 & 6.19 \\
L2 history & 0.294 & rolling mean & 0.012 & 13.73 \\
L5 interaction & 0.164 & score \(\times\) length & 0.012 & 15.03 \\
L4 embedding & 0.095 & \(e_0\) & 0.012 & 27.13 \\
L1 graph & 0.094 & sibling count & 0.074 & 0.88 \\
L6 context & 0.000 & beam & 1.000 & 0.00 \\
\bottomrule
\end{tabular}
\end{center}

The path still has three committed steps, so beam, maximum depth, and the last-step fraction do not move.
Sibling count does move, by \(0.88\), and its concept \(p\)-value is \(0.074\).
That is the graph signal: the drifted judge expands a different frontier.
It is smaller than the change in the score profile along the path.
Score entropy, the rolling mean, and the state encoding all clear \(p=0.012\).
The large embedding shift is the committed partial expression, scaled by 24, landing on different magnitudes after the drift.

\section{One-step offer search}
Beam 1, depth 1, three offers, drift at user 16, judge mix \(0.80\).
\(Y\) moves from \(0.50\) to \(0.167\).
Every user scores exactly three children and commits one step, so the graph columns and the context columns are constant.

\begin{center}
\begin{tabular}{lrrrl}
\toprule
Layer & SSE share & Best column & Best \(p\) & Largest mean shift \\
\midrule
L2 history & 0.732 & previous \(Y\) & 0.012 & 0.333 \\
L5 interaction & 0.136 & depth \(\times\) gap & 0.370 & 0.278 \\
L4 embedding & 0.132 & \(e_2\) & 0.420 & 1.000 \\
L1 graph & 0.000 & depth & 1.000 & 0.000 \\
L3 statistics & 0.000 & score std & 1.000 & 0.000 \\
L6 context & 0.000 & beam & 1.000 & 0.000 \\
\bottomrule
\end{tabular}
\end{center}

Previous \(Y\) is the column that carries the level shift: once failures cluster, the last outcome predicts the next one, with concept \(p=0.012\) and a mean shift of \(0.333\).
The offer identity is a four-bin encoding. Its mean shifts by \(1\), which is the habit-to-click change, and its concept \(p\)-value stays \(0.42\).
The encoding is a relabeling of the batch bit, so the interaction adds almost no SSE beyond \(Y\sim 1+\mathrm{batch}\).
A single score per episode has standard deviation zero, and the statistics layer is empty for the same reason.

\section{What the attribution supports}
On a search whose frontier changes, FSDS assigns the outcome drop first to the score profile, then to the state encoding, and only then to sibling count.
On a search whose tree shape is fixed, the graph layer is silent, the identity encoding is visible as a mean shift and not as a concept column, and the usable concept column is the previous outcome.
Context columns that are constant by design contribute share zero in both tasks.
No budget weight was imposed on the layers. The shares above are the SSE sums.

\end{document}
